\section{Further theorems of Fredholm type}

1
2 Is this something already known from finite-dimensional linear algebra?
3
4 
5 
6
7 In finite dimensions, this is where the characteristic polynomial eventually comes from
8
9
10
11
12
13 Is this something already known from finite-dimensional linear algebra?
14 What does this look like in a Hilbert space and from what theorem of chapter 3 does this follow?
15 So lambda=0 is also a weird case here

\begin{exercise}{1}
Show that $f_0(y) = f_0(T_\lambda x) = g(x)$, for a given $g \in X'$, in the proof of Theorem 8.5-3 is linear.
\end{exercise}
\begin{proof}
This is direct from the linearity of $g$.
\end{proof} 

\begin{exercise}{2}
What does Theorem 8.5-1 imply in the case of a system of $n$ linear algebraic equations in $n$ unknowns?
\end{exercise}
\begin{proof}
fill
\end{proof} 

\begin{exercise}{3}
Consider a system $Ax - y$ consisting of $n$ linear equations in $n$ unknowns.
Assuming that the system  has a solution $x$, show that $y$ must satisfy a condition of the form $f(y) = 0$ for all $f \in X'$ satisfying $T^\times f - \lambda f = 0$.
\end{exercise}
\begin{proof}
fill
\end{proof} 

\begin{exercise}{4}
A system $Ax = y$ of $n$ linear equations in $n$ unknowns has a (unique) solution for any given $y$ if and only if $Ax = 0$ has only the trivial solution $x = 0$.
How does this follow from one of our present theorems?
\end{exercise}
\begin{proof}
fill
\end{proof} 

\begin{exercise}{5}
fill
\end{exercise}
\begin{proof}
fill
\end{proof} 

\begin{exercise}{6}
fill
\end{exercise}
\begin{proof}
fill
\end{proof} 

\begin{exercise}{7}
fill
\end{exercise}
\begin{proof}
fill
\end{proof} 

\begin{exercise}{8}
fill
\end{exercise}
\begin{proof}
fill
\end{proof} 

\begin{exercise}{9}
fill
\end{exercise}
\begin{proof}
fill
\end{proof} 

\begin{exercise}{10}
fill
\end{exercise}
\begin{proof}
fill
\end{proof} 

\begin{exercise}{11}
fill
\end{exercise}
\begin{proof}
fill
\end{proof} 

\begin{exercise}{12}
fill
\end{exercise}
\begin{proof}
fill
\end{proof} 

\begin{exercise}{13}
fill
\end{exercise}
\begin{proof}
fill
\end{proof} 

\begin{exercise}{14}
fill
\end{exercise}
\begin{proof}
fill
\end{proof} 

\begin{exercise}{15}
fill
\end{exercise}
\begin{proof}
fill
\end{proof} 
